\documentclass[11pt,a4paper]{article}
\usepackage[margin=27mm]{geometry}
\usepackage{amsmath,amsthm,amssymb,fontspec,unicode-math,hyperref}
\usepackage{longtable,booktabs,array,calc,fvextra}
\setmainfont{STIXGeneral}
\setmathfont{STIX Two Math}
\setmonofont{Menlo}
\hypersetup{hidelinks}
\newtheorem{theorem}{Theorem}
\newtheorem{proposition}[theorem]{Proposition}
\newtheorem{lemma}[theorem]{Lemma}
\newtheorem{corollary}[theorem]{Corollary}
\theoremstyle{remark}\newtheorem{remark}[theorem]{Remark}
\newcommand{\gen}[1]{\langle #1\rangle}
\newcommand{\F}{\mathbb F}
\newcommand{\Z}{\mathbb Z}
\newcommand{\Aut}{\operatorname{Aut}}
\newcommand{\Sym}{\operatorname{Sym}}
\setlength{\emergencystretch}{2em}
\providecommand{\tightlist}{\setlength{\itemsep}{0pt}\setlength{\parskip}{0pt}}

\title{Regular polyhedra with equal adjacent parameters}
\author{Human Compression Project}
\date{}
\begin{document}
\maketitle
\begin{abstract}
We recover the exact matrix model and universal presentation of
Li--Zhang for a self-dual regular polyhedron of type $\{n,n\}$ and
order $n2^n$, for every $n\ge3$. The group is the augmentation module
of an $n$-cycle extended by its dihedral symmetry group. The common
theorem recognizing the translation subgroup supplies an explicit inverse from this
model to the presentation. We also identify the distinguished marking
with the earlier polygon instance of Montero--Ivi\'c Weiss's halving
construction, and state the parameter specializations, the exceptional
value $n=2$, and the necessary correction to a displayed multiplication
convention.
\end{abstract}
\section{Statement and provenance}
Write $D_{2n}$ for the dihedral group of order $2n$ and
$E_n=\{v\in\F_2^{\Z/n\Z}:\sum_i v_i=0\}$. We use the common paper
[U], specifically its graph generation, complement twist and cyclic
recognition results. Their hypotheses are verified below. We retain the
standard correspondence between string C-groups and regular abstract
polytopes [MS02, Section 2E].

\begin{theorem}\label{C4.main}
For each integer $n\ge3$, the presentation
\[
\begin{split}
P_n=\langle r_0,r_1,r_2\mid{}&
r_0^2=r_1^2=r_2^2=(r_0r_2)^2=1,\\
&(r_0r_1)^n=(r_1r_2)^n=1,\\
&((r_1r_2)^k(r_0r_1)^k)^2=1\quad(1\le k<n)\rangle
\end{split}
\]
is, with its distinguished generators, the group $E_n\rtimes D_{2n}$
described below. It is a string C-group of type $\{n,n\}$, has order
$n2^n$, and has an involutory automorphism reversing the marking.
Its regular abstract polyhedron is therefore self-dual and has exactly
$n2^n$ flags.
\end{theorem}

This is the principal construction of [LZ25, Theorems 1.3 and 3.4].
The same polyhedron occurs in [MW20, Section 4, physical page 9].
We recover both its marking and the exact later presentation; we do not
claim a new existence result. The higher-rank Conjecture 3.7 in [LZ25]
is a separate question and is not used in this proof.

\section{The marking and its geometry}
Let $H=\gen{a,b\mid a^n=b^2=1,\ bab=a^{-1}}$ act on
$\Omega=\Z/n\Z$ by $a(i)=i+1$, $b(i)=-i$.
Put $s=ab$ and $q=e_0+e_1$. In the left semidirect product
\[
 G_n=E_n\rtimes H,\qquad (v,g)(w,h)=(v+gw,gh),
\]
take
\begin{equation}\label{C4.marking}
 r_0=(0,s),\qquad r_1=(0,b),\qquad r_2=(q,s).
\end{equation}
Since $sq=q$, these are nonidentity involutions. Also
$r_0r_2=(q,1)$ and $r_0r_1=(0,a)$.
They generate $H$ and one edge translation of the cycle, so the common
graph-generation corollary gives the whole group $G_n$.

Apply the common complement-twist lemma with base point $0$:
\[
 \delta(v,h)=(v+e_0+he_0,h).
\]
It fixes $r_1$ and interchanges $r_0,r_2$. Thus both adjacent products
have order exactly $n$. The two dihedral subgroups
\[
 A=\gen{r_0,r_1}=\{(0,h):h\in H\},\qquad
 B=\gen{r_1,r_2}=\delta(A)
\]
intersect in the stabilizer of $0$, namely $\gen{r_1}$.
The remaining maximal marked subgroup is
$C=\gen{r_0,r_2}=\{1,r_0,r_2,(q,1)\}$.
Inspecting its translation coordinates gives $A\cap C=\gen{r_0}$;
applying $\delta$ gives $B\cap C=\gen{r_2}$.

Each of $A,B,C$ is a rank-two C-group. If two proper marked subgroups
lie in the same one, its intersection property applies. Otherwise
choose two different maximal marked subgroups containing them.
Their intersection is the subgroup generated by their common marker,
as just proved. Intersecting with that subgroup inside each maximal
group gives exactly the subgroup on the common index set.
Intersections with the full group are tautological. This proves the
full intersection property for the three generators. Their distant
pair commutes, so the marking is a string C-group. The retained
correspondence now gives the polyhedron and its full flag-regular
group. The reversal $\delta$ supplies its duality.

The basis $e_0+e_i$, $i\ne0$, gives $|E_n|=2^{n-1}$, hence
$|G_n|=n2^n$. This count describes the group already identified;
it will not be used to prove the universal presentation.

\section{An inverse map for the presentation}
In $P_n$ put $x=r_0r_1$, $y=r_1r_2$ and $b=r_1$.
The defining relations give
\[
 x^n=y^n=b^2=1,\quad bxb=x^{-1},\quad byb=y^{-1},
 \quad(y^kx^k)^2=1\quad(1\le k<n).
\]
In the model~\eqref{C4.marking}, the first two images are $(0,a)$,
$(0,b)$, while
\[
 y=(e_0+e_{-1},a^{-1})=\delta(0,a^{-1}).
\]
Consequently $y^kx^k=(e_0+e_{-k},1)$ for every integer $k$.
All relators of $P_n$ hold in $G_n$, including $(r_0r_2)^2=1$,
so the marking gives a homomorphism $\Phi:P_n\to G_n$.

The common cyclic-recognition corollary applies to these exact
relations. It constructs the inverse by
\[
 (e_0+e_i,1)\longmapsto y^{-i}x^{-i},\qquad
 (0,a)\longmapsto x,\qquad(0,b)\longmapsto b.
\]
It is well defined on indices modulo $n$, and the common proof verifies
the conjugation relations needed for a semidirect-product homomorphism.
The composites fix all translation basis vectors and dihedral
generators, and all generators of $P_n$, respectively.
Thus $\Phi$ is an isomorphism. Together with the preceding section,
this proves Theorem~\ref{C4.main} without a finite-order upper bound.

\section{Exact identifications and parameters}
The source uses the ordered basis
\[
 f_j=e_0+e_{n-j}\qquad(1\le j<n)
\]
of $E_n$. Direct application of the two permutations gives
\[
 bf_j=f_{n-j},\qquad af_1=f_{n-1},\qquad
 af_j=f_{j-1}+f_{n-1}\quad(2\le j<n).
\]
These are exactly its matrices $V,A$; the matrix of $s=ab$ is $U=AV$,
and its last coordinate vector is $f_{n-1}=q$. This verifies the literal
matrix marking in [LZ25, Section 3], in the consistent left-action
convention used here.

For the earlier geometric construction, use the standard marking
$(e_0,1),(0,s),(0,b)$ of the generalized cube over the $n$-gon in
$\F_2^\Omega\rtimes H$. Endpoint halving changes the first marker to
\[
 (e_0,1)(0,s)(e_0,1)=(q,s)
\]
and leaves $(0,s),(0,b)$. Ordering the new three-node diagram as a
string and reversing it gives~\eqref{C4.marking}. Thus this is the
same distinguished group, not merely another group of the same order.
The halving theorem and polygon count are proved in the common paper
and the companion application to generalized cubes; the direct
intersection proof above also establishes this application's geometric
conclusion independently of that construction theorem.

Every odd $n\ge3$ gives [LZ25, Corollary 3.5]. More generally,
putting $n=2^sp\ge3$, for integers $p\ge2$, $s\ge0$, gives
\[
 |G_n|=2^{2^sp+s}p.
\]
This includes the prime range of its Corollary 3.6 and also all composite
$p\ge2$. The value $(p,s)=(2,0)$ requires a separate construction:
$C_2^3$ with its three coordinate involutions is a string C-group of
type $\{2,2\}$, has eight flags and admits reversal of the coordinates.
No faithful permutation action of $D_4$ on two vertices is assumed.

\section{Conventions, attribution and formal scope}
The multiplication displayed in [LZ25, physical page 7] uses
$\varphi(g)^{-1}w$ while describing $\varphi$ as a homomorphism acting
on column vectors. With left composition this is not generally
associative. For $n=3$, take
$u=(0,a)$, $v=(0,b)$ and $w=(e_0+e_1,1)$.
That rule gives translation coordinate $e_0+e_1$ in $(uv)w$ and
$e_1+e_2$ in $u(vw)$. Our multiplication uses $\varphi(g)w$.
A right-action convention is also possible, but its composition order
must change consistently. The recovered principal statements do not
assert the literal validity of that source display.

The original existence and presentation results remain attributed to
the two cited papers. The common displacement-commutativity proof
expands the argument of Li--Zhang. The present proof change is the use
of the common augmentation-module recognition theorem to construct
the inverse explicitly; the previous normal-subgroup order bound is
no longer required. No priority claim is made for that general method.

The existing file \texttt{BinaryDihedralDisplacements.lean} verifies
the integer displacement identity, its commutative closure and the
finite-period adapter, with five reported declarations using only
standard axioms. The companion file \texttt{SemidirectRecognition.lean} proves the
generic inverse construction from actual component homomorphisms,
covariance and generation, with nine reported declarations. The
augmentation basis, its cyclic instantiation here, intersection property,
matrix identification and geometric correspondence remain written proofs.
Neither collection is a complete formalization of
Theorem~\ref{C4.main}.
The former finite model program is retained as an optional diagnostic,
not a premise for the quantified theorem.

\section*{References}
{[}U{]} \emph{Changes of chamber generators and binary translation groups},
the common paper of Topic C.

{[}MW20{]} A. Montero and A. Ivi\'c Weiss,
\emph{Locally spherical hypertopes from generalised cubes},
arXiv:2001.11082v1, 29 January 2020.

{[}LZ25{]} M. Li and W.-J. Zhang,
\emph{Regular 3-polytopes of type $\{n,n\}$},
arXiv:2505.09505v1, 14 May 2025. Theorems 1.3 and 3.4,
Corollaries 3.5--3.6.

{[}MS02{]} P. McMullen and E. Schulte,
\emph{Abstract Regular Polytopes}, Cambridge University Press, 2002,
Section 2E.
\end{document}
